% Classe do documento - PREÂMBULO
\documentclass[a4paper, oneside, 11pt,]{report} %showframe
%\documentclass[chapterprefix=false,appendixprefix=true]{scrreprt}
\usepackage[printonlyused,withpage]{acronym}
\usepackage[T1]{fontenc}
\usepackage[utf8]{inputenc} %%%caracteres latinos
%\usepackage{lmodern}
\usepackage{txfonts}
\usepackage[a4paper, left=3cm, right=2cm, top=4.3cm, bottom=3.7cm]{geometry}
%%%%%%%%%%%%%%%%%
%\usepackage[printonlyused,withpage]{acronym}
%%%%%%%%%%%%%%%%%%%%%%%%%
\usepackage{enumitem}
\usepackage{wasysym}
\usepackage{blindtext}
\usepackage{multirow}
\usepackage{rotating}
\usepackage{slashbox}
\usepackage{booktabs}
\usepackage{tabularx}
\usepackage{longtable}
%\usepackage{ctable}
%\usepackage{color}
%\usepackage{europs}%simbolo euro
%\usepackage[gen]{eurosym}
\usepackage{placeins} 
\usepackage{savesym}
\savesymbol{iint}
\usepackage{txfonts}
\restoresymbol{TXF}{iint}
% c o m m a n d  \text  f r o m   a m s m a t h
%%%%%%%%%%%%%%%%%%%%%%%%%%%%%%%%%%%%%%%%%%%
\makeatletter
\DeclareRobustCommand{\text}{%
	\ifmmode\expandafter\text@\else\expandafter\mbox\fi}
\let\nfss@text\text
\def\text@#1{{\mathchoice
		{\textdef@\displaystyle\f@size{#1}}%
		{\textdef@\textstyle\f@size{#1}}%
		{\textdef@\textstyle\sf@size{#1}}%
		{\textdef@\textstyle \ssf@size{#1}}%
		\check@mathfonts
	}%
}
\def\textdef@#1#2#3{\hbox{{%
			\everymath{#1}%
			\let\f@size#2\selectfont
			#3}}}
\makeatother
% c o m m a n d  \text  f r o m   a m s m a t h
%%%%%%%%%%%%%%%%%%%%%%%%%%%%%%%%%%%%%%%%%%%

%\usepackage{helvet}

\newtheorem{theorem}{Proposição}

\usepackage{pgfplots}
\pgfplotsset{compat=newest}
%\usepackage{framed}
%\input{nicetoc}
\usepackage{etoolbox}
\usepackage{pstricks-add}


\pgfdeclarelayer{background}
\pgfdeclarelayer{foreground}
\pgfsetlayers{background,main,foreground}



%\usepackage{floatrow} % Table float box with bottom caption, box width adjusted to content
%	\floatsetup[table]{style=plaintop}
%	\DeclareFloatSeparators{MySep}{\hspace{-12pt}}
%	\floatsetup[figure]{style=plain, floatrowsep=Mysep}
%	\floatsetup[subfigure]{subfloatrowsep=none}

\usepackage{marvosym}
\usepackage{import}
\usepackage[normalem]{ulem} %permite sublinhar palavras

\usepackage[activate={true,nocompatibility}, final, tracking=true, kerning=true, spacing=true factor=1100,stretch=10,shrink=10]{microtype}
\SetProtrusion{encoding={*},family={bch},series={*},size={6,7}}
{1={ ,750},2={ ,500},3={ ,500},4={ ,500},5={ ,500},
	
	6={ ,500},7={ ,600},8={ ,500},9={ ,500},0={ ,500}}
\microtypecontext{spacing=nonfrench}

\SetExtraKerning[unit=space]
{encoding={*}, family={bch}, series={*}, size={footnotesize,small,normalsize}}
{\textendash={400,400}, % en-dash, add more space around it
	"28={ ,150}, % left bracket, add space from right
	"29={150, }, % right bracket, add space from left
	\textquotedblleft={ ,150}, % left quotation mark, space from right
	\textquotedblright={150, }} % right quotation mark, space from left

\SetExtraKerning[unit=space]
{encoding={*}, family={qhv}, series={b}, size={large,Large}}
{1={-200,-200}, 
	\textendash={400,400}}

\SetTracking{encoding={*}, shape=sc}{40}

% activate={true,nocompatibility} - activate protrusion and expansion
% final - enable microtype; use "draft" to disable
% tracking=true, kerning=true, spacing=true - activate these techniques
% factor=1100 - add 10% to the protrusion amount (default is 1000)
% stretch=10, shrink=10 - reduce stretchability/shrinkability (default is 20/20)

%\usepackage{datetime}
\usepackage{etoolbox}
\usepackage[flushleft]{threeparttable}

\usepackage{lscape}
\usepackage{multicol}
% -> tabela colorida
\usepackage{colortbl}
\usepackage{array}
\newcolumntype{L}[1]{>{\raggedright\let\newline\\\arraybackslash\hspace{0pt}}m{#1}}
\newcolumntype{C}[1]{>{\centering\let\newline\\\arraybackslash\hspace{0pt}}m{#1}}
\newcolumntype{R}[1]{>{\raggedleft\let\newline\\\arraybackslash\hspace{0pt}}m{#1}}

%\usepackage[table]{xcolor}
%%%%%%%%%%%%%%
\usepackage{pdfpages}
\usepackage{wrapfig}
%\usepackage{cutwin}
%%%%%%%%%%%%%%%%%%%%%%%%%%%%%%%%%%%%%%%%%%%%%%%%%%%%%%
\setcounter{secnumdepth}{3} % profundidade do sumário
\setcounter{tocdepth}{3} % % profundidade do sumário
%%%%%%%%%%%%%%%%%%%%%%%%%%%%%%%%%%%%%%%%%%%%%%%%%%
\usepackage{graphicx}
\usepackage{tikz}
\usetikzlibrary{fit,
	arrows,
	calc,
	positioning,
	fillbetween,
	matrix,
	plotmarks,
	dateplot}
\usetikzlibrary{fadings}
\usepackage[skins]{tcolorbox}
\usepackage{pgfplots}
\usepackage{pgfplotstable}


\usepackage{indentfirst}
\usepackage[brazil]{babel}
\usepackage{setspace}
\onehalfspacing

\usepackage{pdfpages}
%%%%%%%%%%%%%%%%%%%%%%%%%%%%%%%%%%%%%%%%%%%%%%%%%%%%%%%%%%%%%%
\usepackage{caption}
\captionsetup{labelsep=endash,
	font=small,
	labelfont=bf,
	textfont=small,
	listformat=simple,
	figurewithin=none,
	tablewithin=none}
%\renewcommand{\thetable}{\thechapter.\arabic{table}}
%\renewcommand{\thefigure}{\arabic{figure}}
\usepackage{subcaption}
\renewcommand{\thesubfigure}{Fig. \thefigure.\arabic{subfigure}}
%	\DeclareCaptionSubType*[labelsep]{subfigure}
\captionsetup[subfigure]{labelformat=simple,labelsep=space}
%	listofformat=subsimple}
%	\captionsetup{lofdepth=2}
%	\makeatletter
%	\renewcommand{\p@subfigure}{}
%\makeatother

%\captionsetup[table]{name = Quadro}
\newcommand{\aref}[1]{\cref{#1} \textendash ~\nameref{#1}} % Anexo A - Título do Anexo
%%%%%%%%%%%%%%%%%%%%%%%%%%%%%%%%%%%%%%%%%%%%%%%%%%%%%%%%
%\thispagestyle{plain}
\usepackage{fancyhdr}
\usepackage{titletoc}
\usepackage{titlesec}

%\setcounter{secnumdepth}{3}

%%% cabeçalho e rodapé p/ páginas
\pagestyle{fancy}
\fancyhead{} % clear all header fields
\chead{T \& C Serviços Geoambientais Ltda.}
\rhead{\thepage}
\fancyfoot{} % clear all footer fields
\cfoot{Catuana Agropecuária LTDA}
\renewcommand{\headrulewidth}{1pt} %espessura da linha do cabeçalho
\renewcommand{\footrulewidth}{1pt} %espessura da linha do rodapé
\fancypagestyle{plain}{ % cabeçalho e rodapé p/ início de capítulos
\fancyhead{} % clear all header fields
\chead{T \& C Serviços Geoambientais Ltda.}
\rhead{\thepage}
\fancyfoot{} % clear all footer fields
\cfoot{Catuana Agropecuária LTDA}
\renewcommand{\headrulewidth}{1pt} %espessura da linha do cabeçalho
\renewcommand{\footrulewidth}{1pt}} %espessura da linha do rodapé
%%%%%%%%%%%%%%%%%%%%%%%%%%%%%%%%%%%%%%%%%%%%%%%%%%%%%%%%%%%%%%%%%%%%%
\setlength{\headheight}{14.5pt}
%% F O R M A T A Ç Ã O  D E  C A P Í T U L O S ,  S E Ç Õ E S ...
%%%%%%%%%%%%%%%%%%%%%%%%%%%%%%%%%%%%%%%%%%%%%%%%%%
\titleclass{\part}{top} % make part like a chapter
\titleformat{\part}
[display]
{\centering\normalfont\Large\bfseries}
{\titlerule[2pt]\vspace{3pt}\titlerule[1pt]\vspace{3pt}\MakeUppercase{\partname} \thepart}
{0pt}
{\titlerule[1pt]\vspace{1pc}\Large\MakeUppercase}
%
\titlespacing*{\part}{0pt}{0pt}{20pt}
%
%%%%%%%%%%%%%%%%%%%%%%%%%%%%%%%%%%%%%%%%%%%%%%%%%%
\titlespacing*{\chapter}{0pt}{10pt}{20pt} % espaço entre o cabeçalho e capítulo
%
\titleformat{\chapter}[display]
{\normalfont %\sffamily
\large \bfseries}
{\chaptertitlename\ \thechapter}{16pt}{\large \MakeUppercase }
\titleformat{\section}
{\normalfont %\sffamily
\normalsize \bfseries}
{\thesection}{1em}{}
\titleformat{\subsection}
%% subsection format
{\normalfont %\sffamily
\normalsize \bfseries}
{\thesubsection}{1em}{}
%% subsubsection format
\titleformat{\subsubsection}
{\normalfont \MakeLowercase %\sffamily
\normalsize \bfseries}
{\thesubsubsection}{1em}{}
%% paragraph format
\titleformat{\paragraph}
{\MakeLowercase \normalfont\normalsize\bfseries}{\theparagraph}{1em}{}
\titlespacing*{\subparagraph}{\parindent}{3.25ex plus 1ex minus .2ex}{.75ex plus .1ex}
%% subparagraph format
\titleformat{\subparagraph}
{\MakeLowercase\normalfont\normalsize\bfseries}{\thesubparagraph}{1em}{}
\titlespacing*{\subparagraph}{\parindent}{3.25ex plus 1ex minus .2ex}{.75ex plus .1ex}

%%%% F O R M A T A Ç Ã O  T Í T U L O  D O  S U M Á R I O
\usepackage{tocloft} %change table of contents

\setlength{\cftbeforetoctitleskip}{10pt} % espaço vertical antes do titulo toc
\setlength{\cftaftertoctitleskip}{20pt} % espaço vertical depois do titulo toc

\setlength{\cftbeforeloftitleskip}{10pt} % espaço vertical antes do titulo lof
\setlength{\cftafterloftitleskip}{20pt} % espaço vertical antes do titulo lof

\setlength{\cftbeforelottitleskip}{10pt}% espaço vertical antes do titulo lot
\setlength{\cftafterlottitleskip}{20pt}% espaço vertical antes do titulo lot

%\setlength{\cha}

%add o label Figura - na Lista de Figuras
\newlength{\mylen}
\renewcommand{\cftfigpresnum}{\figurename\enspace}
\renewcommand{\cftfigaftersnum}{ \textendash }
\settowidth{\mylen}{\cftfigpresnum}
\addtolength{\cftfignumwidth}{\dimexpr\mylen}
\setlength{\cftfigindent}{0pt}% no indent

%add o label Tabela - na Lista de Tabelas
\renewcommand{\cfttabpresnum}{\tablename\enspace}
\renewcommand{\cfttabaftersnum}{ \textendash }
\settowidth{\mylen}{0.6\cfttabpresnum\cfttabaftersnum}
\addtolength{\cfttabnumwidth}{0.7\mylen}
\setlength{\cfttabindent}{0pt}% no indent
%sumário upper case

% centralizar sumário
\renewcommand{\cfttoctitlefont}{\hfill \large \bfseries \MakeUppercase} 
\renewcommand{\cftaftertoctitle}{\hfill}

% centralizar lista de figuras
\renewcommand{\cftloftitlefont}{\hfill \large \bfseries \MakeUppercase }
\renewcommand{\cftafterloftitle}{\hfill}

% centralizar lista de tabelas
\renewcommand{\cftlottitlefont}{\hfill \large\bfseries \MakeUppercase }
\renewcommand{\cftafterlottitle}{\hfill}


% % % A N E X O S % % % 

\usepackage{appendix}

\let\oldappendices\appendices
\def\appendices{\oldappendices\adjustmtc}

\appendixpageoff

\newlength{\tocnumwidth}
\newlength{\temp}
\setlength{\tocnumwidth}{\cftchapnumwidth} % store the starting width of the TOC chapter number

%\renewcommand{\cftchappresnum}{\chaptername\space}
%\settowidth{\temp}{\bfseries\cftchappresnum\cftchapaftersnum} % calculate the extra space
%\addtolength{\cftchapnumwidth}{\temp} % add the extra space
% uncomment the 3 lines above  if you want the word Chapter to precede the number in the TOC

\renewcommand*{\cftchapaftersnum}{\space}

\renewcommand{\appendixname}{ANEXO}

\DeclareRobustCommand{\changetoapp}{%
\renewcommand*{\cftchappresnum}{ANEXO\space}%
\renewcommand{\cftchapaftersnum}{~\textendash}
\settowidth{\temp}{\bfseries\cftchappresnum\cftchapaftersnum}
% comment the line above out if you include Chapter before the number in the TOC want the Chapter and Appendix TOC entries to line up.
\renewcommand{\cftsecaftersnum}{~\textendash}

\setlength\cftchapnumwidth{\tocnumwidth}
\addtolength{\cftchapnumwidth}{\temp}
%%
\setlength\cftsecnumwidth{\tocnumwidth}
\addtolength{\cftsecnumwidth}{0.35\temp}}
%%
\let\oldappendix\appendix
\renewcommand*{\appendix}{\oldappendix \addtocontents{toc}{\changetoapp}}

%


%new command to insert a blank page
\newcommand{\insertblankpage}{
\vfill
\newpage
\thispagestyle{empty}
\mbox{}
\newpage
}
%new command to insert a page break
\newcommand{\startnewpage}{
\vfill
\pagebreak[4]
}

%\usepackage{xparse}
%\DeclareDocumentCommand \sectionlevel{ s m m }% s = an optional asterisk, m = mandatory argument
%	{ 
%	\IfBooleanTF{#1}
%	{
	%		\ifnumcomp{#2}{=}{0}{\addcontentsline{toc}{part}{#3} \part*{#3}}{}
	%		\ifnumcomp{#2}{=}{1}{\addcontentsline{toc}{chapter}{#3} \chapter*{\hfill #3 \hfill}}{}
	%		\ifnumcomp{#2}{=}{2}{\addcontentsline{toc}{section}{#3} \section*{#3}}{}
	%		\ifnumcomp{#2}{=}{3}{\addcontentsline{toc}{subsection}{#3} \subsection*{#3}}{}
	%		\ifnumcomp{#2}{=}{4}{\addcontentsline{toc}{subsubsection}{#3} \subsubsection*{#3}}{}
	%		\ifnumcomp{#2}{=}{5}{\paragraph*{#3}}{}
	%		\ifnumcomp{#2}{=}{6}{\subparagraph*{#3}}{}
	%	}
%	{
	%		\ifnumcomp{#2}{=}{0}{\part{#3}} {}
	%		\ifnumcomp{#2}{=}{1}{\chapter{#3}}{}
	%		\ifnumcomp{#2}{=}{2}{\section{#3}}{}
	%		\ifnumcomp{#2}{=}{3}{\subsection{#3}}{}
	%		\ifnumcomp{#2}{=}{4}{\subsubsection{#3}}{}
	%		\ifnumcomp{#2}{=}{5}{\paragraph{#3}}{}
	%		\ifnumcomp{#2}{=}{6}{\subparagraph{#3}}{}
	%	}
%
%	}

%\usepackage{hyperref}
%\hypersetup{
%	bookmarks=true,         % show bookmarks bar?
%	unicode=false,          % non-Latin characters in Acrobat’s bookmarks
%	pdftoolbar=true,        % show Acrobat’s toolbar?
%	pdfmenubar=true,        % show Acrobat’s menu?
%	pdffitwindow=false,     % window fit to page when opened
%	pdfstartview={FitH},    % fits the width of the page to the window
%	pdftitle={My title},    % title
%	pdfauthor={Author},     % author
%	pdfsubject={Subject},   % subject of the document
%	pdfcreator={Creator},   % creator of the document
%	pdfproducer={Producer}, % producer of the document
%	pdfkeywords={keyword1} {key2} {key3}, % list of keywords
%	pdfnewwindow=true,      % links in new window
%	colorlinks=false,       % false: boxed links; true: colored links
%	linkcolor=red,          % color of internal links (change box color with linkbordercolor)
%	citecolor=green,        % color of links to bibliography
%	filecolor=magenta,      % color of file links
%	urlcolor=cyan           % color of external links
%	hidelinks = true
%}


\usepackage{etoc}
%\newcommand*{\figuretotoc}{\etoctoccontentsline{Figure}{\thefig}}



\usepackage{cleveref} %referencia
\usepackage{nameref}
\crefname{part}{Parte}{Partes}
\crefname{chapter}{Capítulo}{Capítulos}
\crefname{section}{Seção}{Seções}
\crefname{subsection}{Seção}{Seções}
\crefname{subsubsection}{Seção}{Seções}
\crefname{figure}{Figura}{Figuras} % changes default behavior to Figure. 1
\crefname{table}{Tabela}{Tabelas}
\crefname{equation}{Equação}{Equações}
\crefname{subfigure}{Figura}{Figuras}
\crefname{subsubfigure}{Figura}{Figuras}
\crefname{appendix}{Anexo}{Anexos}

\makeatletter
\let\titlesec@part\part
\renewcommand{\part}{\@ifstar\part@star\part@nostar}
\def\part@star#1{\NR@gettitle{#1}\titlesec@part*{#1}}
\def\part@nostar{\@ifnextchar[\part@nostar@opt\part@nostar@nopt}
\def\part@nostar@nopt#1{\NR@gettitle{#1}\titlesec@part{#1}}
\def\part@nostar@opt[#1]#2{\NR@gettitle{#2}\titlesec@part[#1]{#2}}
\makeatother

%bordas na imagem
%%%%%%%%%%%%%%%%%%%%%%%%%%%%%%%%%%%%%%%%%%%%%%%%
\setlength\fboxsep{0pt}    %%%%%%%%%%%%%%%%%%
\setlength\fboxrule{0.75pt}%%%%%%%%%%%%%%%%%%
%%%%%%%%%%%%%%%%%%%%%%%%%%%%%%%%%%%%%%%%%%%%%%%%
\begin{document}


%\pagenumbering{roman}  numeração romana
%\fontfamily{phv} \selectfont %%fonte helvetica
\newgeometry{vmargin={3cm,2cm},hmargin={2cm,2cm}}
%  C A P A
\thispagestyle{empty}
\chead{}
\fancypagestyle{plain}{\chead[LE,RO]{}}
\begin{figure}[t!]
\centering
\includegraphics[width = 7 cm]{./imagens/logo.pdf}
\end{figure}
\mbox{}
 
\vspace{6.5cm}

\begin{center}
{\Large \bfseries
Silvia Alcaldes Gomes\\
Laudo Geológico Geotécnico
}
\end{center}
\mbox{}

\vfill

\begin{flushright}
\begin{minipage}[h]{8cm}
{\footnotesize \textcolor{white}{Relatório de sondagem geológica do tipo SPT, conforme diretrizes da ABNT NBR 6484:2001}}
\end{minipage}
\end{flushright}

\vfill

\begin{center}
{\large
{MARINGÁ/PR\\
DEZEMBRO/2021}
}
\end{center}
\newpage

% F A L S A  F O L H A  D E  R O S T O
\thispagestyle{empty}
\begin{figure}[t!]
\centering
\includegraphics[width = 7 cm]{./imagens/logo2.pdf}
\end{figure}
\mbox{}
 
\vspace{6.5cm}
\begin{center}
{\Large \bfseries
Silvia Alcaldes Gomes\\
Laudo Geológico Geotécnico
}
\end{center}

\vfill

\begin{flushright}
\begin{minipage}[h]{8cm}
{\footnotesize \textcolor{white}{Laudo Geológico Geotécnico }}
\end{minipage}
\end{flushright}

\vfill

\begin{center}
{\large \textcolor{white} {
{MARINGÁ/PR\\
DEZEMBRO/2021}}
}
\end{center}
\newpage

% F O L H A  D E  R O S T O
\pagenumbering{arabic}
\thispagestyle{empty}
\begin{figure}[t!]
\centering
\includegraphics[width = 7 cm]{./imagens/logo.pdf}
\end{figure}
\mbox{} 
\vspace{6.5cm}

\begin{center}
{\Large \bfseries
Silvia Alcaldes Gomes\\
Laudo Geológico Geotécnico
}
\end{center}

\vfill

\begin{flushright}
\begin{minipage}[h]{8cm}
{\footnotesize{Laudo Geológico Geotécnico }}
\end{minipage}
\end{flushright}

\vfill

\begin{center}
{\large 
{MARINGÁ/PR\\
DEZEMBRO/2021}
}
\end{center}

\clearpage
% I N S E R I R  P Á G I N A  E M  B R A N C O

\thispagestyle{empty}


\textbf{CONTRATANTE} \\


\begin{tabular}{p{4cm} p{11cm}}
Nome: & Silvia Alcaldes Gomes\\
CPF: & 020.628.409-85\\
End.: &Rua Olímpico, 81 - Jardim Aurora\\
CEP:  & 87.070-240\\
Cidade/UF:  & Maringá/PR\\
%Telefone: & (44) 9 8803 4567\\
\end{tabular}
\vfill

\textbf{EMPRESA RESPONSÁVEL PELA ELABORAÇÃO}\\


\begin{tabular}{p{4cm} p{11cm}}
%\multicolumn{2}{l}{\includegraphics[width=1cm]{./imagens/logo.pdf}}\\
Razão Social: & T \& C Serviços Geoambientais Ltda\\
CNPJ: & 13.257.061/0001-80\\ %13257061000180
End.: & Rua José Rodrigues de Oliveira, 454 - Jardim Monte Cristo\\
CEP: & 87.706-448\\
Município: & Paranavaí/PR  \\
Tel: & (44) 3423 9956  \\
Site: &  www.araucariaengenharia.com.br\\
e-mail: & araucariapa@hotmail.com \\
\end{tabular}
\vfill
%\newpage
%\thispagestyle{empty}

%\textbf{EQUIPE TÉCNICA}\\

\begin{minipage}[b]{0.49\textwidth}
	
	
	\includegraphics[width=\textwidth]{./imagens/assinatura digital.png}
	\vspace{-40pt}
	
	\rule{8cm}{1pt}
	
	Crislaine Aline Campana 
	
	Geóloga
	
	CREA\textendash~PR 103920/D\\
	
	
\end{minipage}

%\begin{tabular}{l p{8cm}}
%
%\textbf{Responsável Técnico:}& \\
%%\midrule
%&  
%%Rosemary Tuzi Domiciliano & Engenheira Civil \textendash~CREA-PR 117336/D\\
%%Juraci Couto Casula & Tecnóloga em Gestão Ambiental \textendash~CRQ-PR 09202411\\
%%\bottomrule
%\end{tabular}
%\vfill

%\textbf{Equipe de Apoio}


%\begin{tabular}{p{6cm} p{8cm}}

%Especialista em Meio Ambiente 
%e Recursos Naturais -  \\

%Rogério Penteado de Souza & Engenheiro Ambiental \\

%\end{tabular}
\vfill
\thispagestyle{empty}


\restoregeometry

\pagenumbering{roman}
\chead{\thepage}
\fancypagestyle{plain}{\chead{T \& C Serviços Geoambientais Ltda.}\rhead{\thepage}}
%\fontfamily{phv} \selectfont %%fonte helvetica



\import{./01-informacoes/}{00}


\appendix

%\addtocontents{toc}{\protect\setcounter{tocdepth}{-1}}

\renewcommand{\appendixtocname}{ANEXOS}%renomeia de apendice p/ anexo (no sumário)
%\noappendicestocpagenum %nao aparece o numero da pagina no sumário
\addtocontents{toc}{\cftpagenumbersoff{chapter}}
\addtocontents{toc}{\cftpagenumbersoff{section}} % Retira a numeração da pagina no sumário
\addappheadtotoc % adiciona a seção ANEXOS no sumário
\renewcommand{\appendixname}{ANEXO}% renomeia o prefixo Apêndice p/ ANEXO
\renewcommand{\appendixpagename}{ANEXOS}% renomeia o prefixo Apêndice p/ ANEXO
%\addtocontents{toc}{\cftpagenumbersoff{chapter}} % Retira a numeração da pagina no sumário


%\appendixpage
%\addcontentsline{toc}{chapter}{ANEXOS} 
%\chapter*{\hfill {\bfseries ANEXOS} \hfill}
\addtocontents{toc}
{
\setlength{\cftbeforechapskip}{\cftbeforesecskip}
\setlength{\cftchapindent}{\cftsecindent}

\setlength{\cftbeforesecskip}{\cftbeforesubsecskip}
\setlength{\cftsecindent}{\cftsubsecindent}

%\protect\renewcommand{\cftchapfont}{\cftsecfont}
%\protect\renewcommand{\protect\cftchapdotsep}{\cftsecdotsep}
}

\titleclass{\part}{top} % make part like a chapter
\titleformat{\part}
[display]
{\centering\normalfont\large\bfseries}
{\vspace{3pt}\vspace{3pt}\MakeUppercase{\partname} \thepart}
{0pt}
{\titlerule[0pt]\vspace{1pc}\large\MakeUppercase}
%
\titlespacing*{\part}{0pt}{10pt}{20pt}

%\insertblankpage

\part*{ANEXOS}

\etoctoccontentsline*{appendixplaceholder}{APPENDIX}{-1}
\etocsettocdepth{chapter}

\etocmulticolstyle[1]{\normalsize
\hrulefill
\MakeUppercase{}
}
\localtableofcontents
\vspace{-11pt}
\hrulefill

\renewcommand{\thechapter}{\Alph{chapter}} % Usar letras ao inves de numeros para os capitulos
\renewcommand{\thesection}{\thechapter.\Roman{section}} % Usar letras ao inves de numeros para os capitulos
\setcounter{chapter}{0} % resetar a contagem

\clearpage

\fancypagestyle{plain}{\fancyhead[LE,RO]{}}
\pagestyle{fancy}{\fancyhead[LE,RO]{}}

\import{./02-anexos/}{01-lista_anexos}

\end{document}
